\documentclass[a4paper,11pt]{article}

\usepackage[T1]{fontenc}
\usepackage{lmodern}
\usepackage{amsmath}
\usepackage{amssymb}
\usepackage{bm}
\usepackage{graphicx}
\usepackage{booktabs}
\usepackage{tabularx}
\usepackage{array}
\usepackage[section]{placeins}
\usepackage[hidelinks]{hyperref}

\newcolumntype{Y}{>{\raggedright\arraybackslash}X}
\renewcommand{\arraystretch}{1.16}
\setlength{\emergencystretch}{2em}

\title{Learnable Dual-Phase PINNs with Positive- and Negative-Energy
Eigenspace Decomposition for the Time-Dependent Dirac Equation\\
\large Research Overview and Interim Progress Report}
\author{}
\date{September 1, 2026}

\begin{document}

\maketitle

\begin{quote}
\noindent\textbf{Current status.}
The theoretical analysis, single-wavenumber transfer test, method construction,
mechanism ablation, two-wavenumber extension, and continuous-wave-packet
experiment have been completed. The next stage introduces momentum coupling
through an external potential.
\end{quote}

\begin{abstract}
This study investigates long-time phase approximation in physics-informed
neural networks for the time-dependent Dirac equation. For superpositions of
the positive- and negative-energy branches, a single common phase does not
separate both rapid temporal factors with opposite signs. Its numerical benefit
therefore deteriorates under a fixed network size and training budget. A
learnable two-phase branch-aligned network, denoted by LTPB, is introduced to
address this difficulty. The spinor is first projected onto the positive- and
negative-energy eigenspaces of the free Dirac Hamiltonian, after which an
independent real amplitude and learnable phase are assigned to each component.
Single-wavenumber experiments, parameter-matched ablations, and a
two-wavenumber test show that eigenspace alignment and phase separation play
complementary roles. A shared-spectral LTPB subsequently represents 33 Fourier
modes with a fixed parameter count and preserves the correct dispersion phase
at untrained midpoint wavenumbers. Relative to a general Dirac PINN and a
single-common-phase CPA representation, the proposed coordinates may improve
the conditioning of mixed-energy, strongly oscillatory problems by extracting
their rapid phases explicitly. Relative to mature spectral solvers, the
prospective value lies in continuous surrogate construction, sparse-data
assimilation, and parameter inference. These latter advantages remain to be
established through external-potential, noisy-data, and cost-comparison studies.
\end{abstract}

\section{Problem formulation}

Consider the one-dimensional time-dependent Dirac equation
\begin{equation}
i\hbar\partial_t\psi=
\left(-i\hbar c\sigma_x\partial_x+mc^2\sigma_z+V(x,t)I_2\right)\psi,
\qquad
\psi=(\psi_1,\psi_2)^{\mathsf T}.
\end{equation}
Hermiticity of the Dirac Hamiltonian yields the probability density, probability
current density, and continuity equation
\begin{equation}
\rho=\psi^\dagger\psi,
\qquad
j=c\psi^\dagger\sigma_x\psi,
\qquad
\partial_t\rho+\partial_xj=0.
\end{equation}

Suppose that the spinor is parameterized by a common phase,
\begin{equation}
\psi=Ae^{iS/\hbar}.
\end{equation}
Direct substitution gives
\begin{equation}
\rho=A^\dagger A,
\qquad
j=cA^\dagger\sigma_xA.
\end{equation}
The global \(U(1)\) phase \(S\) cancels from both quantities. If the two spinor
components are instead written as
\begin{equation}
\psi_j=a_je^{iS_j/\hbar},
\end{equation}
with real amplitudes \(a_j\), then
\begin{equation}
j=2ca_1a_2\cos\!\left(\frac{S_2-S_1}{\hbar}\right).
\end{equation}
The current density is therefore sensitive to the internal relative phase.
Positive- and negative-energy eigenstates of the free Dirac Hamiltonian carry
the temporal factors \(e^{-iEt}\) and \(e^{+iEt}\), respectively. A single
common phase can absorb at most one of these rapid factors, leaving the other
oscillation to be represented by the complex-amplitude network.

Throughout this report, performance deterioration refers to increased
optimization difficulty and approximation error under a fixed architecture and
training budget. The single-common-phase ansatz retains sufficient formal
representation capacity.

\section{Research program}

\begin{table}[htbp]
\centering
\small
\caption{Research program and current status}
\begin{tabularx}{\textwidth}{@{}c l c Y@{}}
\toprule
Stage & Task & Status & Main result or objective \\
\midrule
1 & Theoretical analysis & Complete &
The common phase does not directly encode the relative phase between the
positive- and negative-energy eigenspaces. \\
2 & Single-mode transfer test & Complete &
CPA-PINN is accurate on a pure positive-energy state, whereas its improvement is
substantially reduced for a mixed-energy state. \\
3 & LTPB construction & Complete &
Independent learnable phases recover the mixed-energy dynamics. \\
4 & Structural ablation & Complete &
Eigenspace alignment and phase separation jointly determine the observed
accuracy. \\
5 & Two-wavenumber test & Complete &
The representation extends from two temporal frequencies to four energy
components. \\
6 & Continuous wave packet & Complete &
A shared-spectral network represents 33 modes and interpolates to untrained
wavenumbers. \\
7 & External potential or scattering & Next &
Test the method under dynamically coupled momentum channels. \\
8 & Sparse noisy data & Planned &
Assess state reconstruction and parameter inference. \\
9 & Full benchmark & Planned &
Compare accuracy, stability, and cost with modern PINNs and classical solvers. \\
10 & Manuscript development & Planned &
Assemble the theory, algorithm, ablations, scope, and data applications into a
complete argument. \\
\bottomrule
\end{tabularx}
\end{table}

\section{Stage I: theoretical analysis}

\subsection{Free positive- and negative-energy branches}

In natural units \(\hbar=c=1\), spatial Fourier transformation reduces the free
Dirac Hamiltonian at wavenumber \(k\) to
\begin{equation}
H(k)=k\sigma_x+m\sigma_z,
\qquad
H(k)u_\pm=\pm E_k u_\pm,
\qquad
E_k=\sqrt{k^2+m^2}.
\end{equation}
As \(k\) varies, the eigenvalues \(+E_k\) and \(-E_k\) form the positive- and
negative-energy branches of the dispersion relation. At fixed \(k\), their
normalized eigenvectors are denoted by \(u_\pm(k)\). A mixed-energy state is
\begin{equation}
\psi_\eta(x,t)=
\frac{e^{ikx}}{\sqrt{L(1+\eta^2)}}
\left(u_+e^{-iE_kt}+\eta u_-e^{+iE_kt}\right).
\end{equation}
For \(\eta=1\), the total norm remains constant while the current density
contains an interference term at frequency \(2E_k\). This single-mode state is
the minimal problem for testing relative-phase approximation.

\subsection{Information carried by conservation constraints}

Let \(R_D=i\partial_t\psi-H_D\psi\) be the Dirac-equation residual. Direct
calculation gives
\begin{equation}
\partial_t\rho+\partial_xj
=2\operatorname{Im}(\psi^\dagger R_D).
\end{equation}
The continuity equation is a real projection of the governing residual along
the state, whereas total norm and total energy are integral constraints. These
conditions restrict the global evolution but do not determine the local
relative phase between the positive- and negative-energy components. This
observation motivates an explicit phase variable for each eigenspace.

\section{Stage II: transfer test of a single-common-phase CPA representation}

\subsection{Experimental setting}

The parameters are \(m=5\), \(k=2\), \(T=5\), and \(L=2\pi\), with periodic
boundary conditions. Two cases are considered:
\begin{itemize}
  \item \texttt{Case P}: a pure positive-energy state;
  \item \texttt{Case M}: an equal-weight superposition of the positive- and
  negative-energy states.
\end{itemize}
Every trainable model uses four hidden layers with 64 \texttt{tanh} neurons per
layer. Adam optimization is run for 5000 steps with 1024 collocation points per
step. Three random seeds are used.

\subsection{Results}

\begin{table}[htbp]
\centering
\footnotesize
\caption{Transfer-test results for the single-common-phase CPA representation}
\begin{tabularx}{\textwidth}{@{}Ycccc@{}}
\toprule
Model &
\shortstack{Case P\\spinor error} &
\shortstack{Case M\\spinor error} &
\shortstack{Case M\\current error} &
\shortstack{Case M\\fidelity} \\
\midrule
Vanilla & 0.9228 & 0.9452 & 0.9693 & 0.2254 \\
SP-CPA & \textbf{0.0067} & 0.7991 & 0.9392 & 0.5935 \\
SP-CPA-Conservative & 0.0111 & 1.0785 & 1.0698 & 0.1631 \\
Analytic two-frequency reference & 0.0079 & \textbf{0.0062} &
\textbf{0.0105} & 0.999990 \\
\bottomrule
\end{tabularx}
\end{table}

\begin{figure}[!htbp]
\centering
\includegraphics[width=0.88\textwidth]{results/figures/01_main_metric_comparison.png}
\caption{The single-common-phase representation is accurate on the pure
positive-energy state, while its improvement is much smaller for the
mixed-energy state.}
\label{fig:cpa-main-en}
\end{figure}

\begin{figure}[!htbp]
\centering
\includegraphics[width=0.88\textwidth]{results/figures/02_case_M_current_x0.png}
\caption{SP-CPA gives a large error for the rapid current oscillation generated
by coherence between the two energy eigenspaces.}
\label{fig:cpa-current-en}
\end{figure}

Relative to Vanilla, SP-CPA reduces the spinor error by \(99.3\%\) in Case P
and by \(15.5\%\) in Case M. Adding continuity, norm, and energy penalties does
not recover the current density or the eigenspace-relative phase. An analytic
wrong-phase construction further shows that nearly correct global invariants
can coexist with substantial local errors in the spinor and current density.

\noindent\textbf{Stage conclusion.}
The common-phase CPA representation remains effective for a pure energy
component. Its rapid-phase extraction is incomplete for a positive/negative
superposition, and its fixed-budget accuracy deteriorates accordingly.

\section{Stage III: LTPB formulation and pilot validation}

\subsection{Method}

Projection onto the eigenspaces of the free Dirac Hamiltonian gives the
learnable two-phase branch-aligned representation
\begin{equation}
\widehat\psi=e^{ikx}
\left[
a_+(x,t)u_+e^{iS_+(x,t)}
+a_-(x,t)u_-e^{iS_-(x,t)}
\right].
\end{equation}
For continuity with the experimental code, this representation is denoted by
LTPB. Here branch alignment is implemented through alignment with the
positive- and negative-energy eigenspaces. The hard initial-condition ansatz is
\begin{equation}
a_\pm=a_{\pm,0}+\tau r_\pm,
\qquad
S_\pm=\tau q_\pm,
\qquad
\tau=t/T.
\end{equation}
Neither the exact energy \(E_k\) nor the exact temporal frequency is supplied
to the forward representation, loss functional, or optimizer. Both phases are
identified from the Dirac residual and the prescribed initial and boundary
conditions.

\subsection{Results}

\begin{table}[htbp]
\centering
\small
\caption{Comparison between SP-CPA and LTPB}
\begin{tabular}{@{}lrrr@{}}
\toprule
Metric & SP-CPA & LTPB & SP-CPA/LTPB \\
\midrule
Case M spinor error & 0.7991 & \textbf{0.0025} & about 320 \\
Current-density error & 0.9392 & \textbf{0.0046} & about 204 \\
Relative-phase RMSE & 1.5864 & \textbf{0.0030} & about 529 \\
Fidelity & 0.5935 & \textbf{0.999996} & -- \\
\bottomrule
\end{tabular}
\end{table}

\begin{figure}[!htbp]
\centering
\includegraphics[width=0.88\textwidth]{results/figures/learned_two_phase_pilot/01_main_metric_comparison.png}
\caption{LTPB attains low errors for both pure- and mixed-energy states without
receiving the exact frequency.}
\label{fig:ltpb-main-en}
\end{figure}

\begin{figure}[!htbp]
\centering
\includegraphics[width=0.88\textwidth]{results/figures/learned_two_phase_pilot/02_learned_branch_phases.png}
\caption{The two phases identified from the governing residual approach
\(-E_kt\) and \(+E_kt\).}
\label{fig:ltpb-phase-en}
\end{figure}

The exact slopes are \(\mp 5.3851648\), while the identified slopes are
\(-5.3855\) and \(+5.3854\). The linear fits have \(R^2\) values close to
unity, and the relative amplitude drift is below \(0.3\%\).

\noindent\textbf{Stage conclusion.}
Independent phase functions assigned to the two energy eigenspaces recover the
relative phase and its associated current oscillation.

\section{Stage IV: structural ablation}

\subsection{Ablation design}

All five models have approximately 12,996 trainable parameters and use paired
initializations and collocation sequences.

\begin{table}[htbp]
\centering
\caption{Structural ablation design}
\begin{tabular}{@{}lcc@{}}
\toprule
Model & Eigenspace aligned & Independent amplitude/phase \\
\midrule
LTPB & Yes & Yes \\
BranchComplex & Yes & No \\
ComponentPolar & No & Yes \\
Rotated LTPB \(15^\circ\) & Partly disrupted & Yes \\
Rotated LTPB \(30^\circ\) & Further disrupted & Yes \\
\bottomrule
\end{tabular}
\end{table}

\subsection{Results}

\begin{table}[htbp]
\centering
\small
\caption{Structural ablation results}
\begin{tabular}{@{}lrrr@{}}
\toprule
Model & Spinor error & Current error & Relative-phase RMSE \\
\midrule
LTPB & \textbf{0.0025} & \textbf{0.0045} & \textbf{0.0029} \\
BranchComplex & 0.9014 & 0.9342 & 1.6268 \\
ComponentPolar & 0.8907 & 0.9309 & 1.4802 \\
Rotated LTPB \(15^\circ\) & 0.8996 & 0.9426 & 1.5579 \\
Rotated LTPB \(30^\circ\) & 0.9306 & 0.9590 & 1.6830 \\
\bottomrule
\end{tabular}
\end{table}

\begin{figure}[!htbp]
\centering
\includegraphics[width=0.88\textwidth]{results/figures/ltpb_mechanism_ablation/01_core_error_comparison.png}
\caption{An eigenspace basis alone and a generic polar representation alone
both remain less accurate than the complete LTPB construction.}
\label{fig:ablation-core-en}
\end{figure}

\begin{figure}[!htbp]
\centering
\includegraphics[width=0.88\textwidth]{results/figures/ltpb_mechanism_ablation/05_rotation_degradation.png}
\caption{Rotation away from the energy eigenspaces rapidly moves the model to a
high-error plateau.}
\label{fig:ablation-rotation-en}
\end{figure}

The preregistered strong criterion required at least a factor of 1.25 error
growth at each additional rotation level. The model had already entered a
high-error plateau at \(15^\circ\), and the increase from \(15^\circ\) to
\(30^\circ\) was only about \(1.02\)--\(1.08\). The strong criterion was
therefore not satisfied. The controlled experiment supports the narrower
conclusion that eigenspace alignment and phase separation are both required for
low error on the present benchmark.

\noindent\textbf{Stage conclusion.}
Parameter matching rules out parameter count and a generic polar
parameterization as sufficient explanations. The observed low error is
associated with the joint eigenspace-aligned, phase-separated construction.
Mechanistic attribution remains limited to the present free single-mode test
because the strong rotation-degradation criterion was not met.

\section{Stage V: two wavenumbers and four energy components}

\subsection{Problem and method}

For \(k_1=2\) and \(k_2=5\), consider
\begin{equation}
\psi=\frac{1}{2\sqrt L}\sum_{l=1}^2
\left[
u_+(k_l)e^{ik_lx-iE_lt}
+u_-(k_l)e^{ik_lx+iE_lt}
\right].
\end{equation}
Multi-LTPB assigns a learnable phase to each of the four \((k_l,\pm)\)
components without using the exact \(E_l\). Five trainable models use nearly
identical parameter counts, 5000 optimization steps, and five paired random
seeds.

\subsection{Results}

\begin{table}[htbp]
\centering
\small
\caption{Two-wavenumber, four-component results}
\begin{tabularx}{\textwidth}{@{}Yrrr@{}}
\toprule
Metric & Multi-LTPB & \shortstack{Best trainable\\baseline} &
\shortstack{Baseline/\\LTPB} \\
\midrule
Spinor error & \textbf{0.006594} & 0.939819 & 142.52 \\
Current-density error & \textbf{0.010976} & 0.957364 & 87.22 \\
Fidelity & \textbf{0.999965} & 0.3172 & -- \\
Paired seeds & \textbf{5/5 pass} & -- & -- \\
\bottomrule
\end{tabularx}
\end{table}

\begin{figure}[!htbp]
\centering
\includegraphics[width=0.88\textwidth]{results/figures/multimode_ltpb/primary_error_comparison.png}
\caption{Multi-LTPB retains low error for two wavenumbers and four energy
components.}
\label{fig:multi-error-en}
\end{figure}

\begin{figure}[!htbp]
\centering
\includegraphics[width=0.88\textwidth]{results/figures/multimode_ltpb/learned_four_phases.png}
\caption{All four projected phase slopes agree with the exact temporal
frequencies; the mean relative slope error is approximately \(0.03\%\).}
\label{fig:multi-phase-en}
\end{figure}

\begin{figure}[!htbp]
\centering
\includegraphics[width=0.88\textwidth]{results/figures/multimode_ltpb/representative_density_current.png}
\caption{The model recovers the density and current interference generated by
multiple spatial modes and both energy eigenspaces.}
\label{fig:multi-current-en}
\end{figure}

\noindent\textbf{Stage conclusion.}
All preregistered accuracy criteria are satisfied. Phase separation extends to
multiple known discrete energy components, while the wavenumbers and
eigenspinor basis remain prescribed from the initial state.

\section{Stage VI: continuous Gaussian wave packet and shared-spectral LTPB}

\subsection{Problem}

On \(L=16\pi\) with \(m=5\) and \(T=5\), 33 wavenumbers
\(k_n=n/8\in[0,4]\) are used. The initial momentum distribution is Gaussian
with center \(k_0=2\) and width \(\sigma_k=0.5\), with equal positive- and
negative-energy weights. The evolution contains dozens of dispersion
frequencies, so assigning separate network outputs to every mode would make the
parameter count grow with spectral resolution.

\subsection{Shared-spectral formulation}

A single network receives \((k,t)\) and returns the two amplitudes and phases:
\begin{equation}
\widehat{\widetilde\psi}(k,t)=
a_+(k,t)u_+(k)e^{iS_+(k,t)}
+a_-(k,t)u_-(k)e^{iS_-(k,t)}.
\end{equation}
Training uses the spinor-form spectral residual
\begin{equation}
R_k=i\partial_t\widetilde\psi
-(k\sigma_x+m\sigma_z)\widetilde\psi.
\end{equation}
The exact dispersion \(E(k)\) is not supplied as supervision. For the free
periodic problem, Parseval's identity makes this residual equivalent to the
corresponding physical-space residual. The parameter count remains 12,932 and
does not scale linearly with the number of modes.

\subsection{Results}

\begin{table}[htbp]
\centering
\footnotesize
\caption{Continuous Gaussian-wave-packet results}
\begin{tabularx}{\textwidth}{@{}Y r r r@{}}
\toprule
Metric & \shortstack{Shared-Spectral-\\LTPB} &
\shortstack{Best trainable\\baseline} & Criterion \\
\midrule
Spinor error & \textbf{0.041038} & 0.909412 & 22.16 \\
Density error & \textbf{0.015262} & 0.487106 & -- \\
Current-density error & \textbf{0.048057} & 0.884010 & 18.40 \\
Fidelity & \textbf{0.998249} & 0.386112 & Pass \\
Paired seeds & \textbf{5/5 pass} & -- & Pass \\
Phase error at training points & \textbf{0.115\%} & -- & \(<1\%\) \\
Phase error at unseen midpoints & \textbf{0.115\%} & -- & \(<2\%\) \\
\bottomrule
\end{tabularx}
\end{table}

\begin{figure}[!htbp]
\centering
\includegraphics[width=0.88\textwidth]{results/figures/wavepacket_shared_spectral_ltpb/primary_error_comparison.png}
\caption{With a fixed parameter count, shared-spectral LTPB is more accurate
than the physical-space and same-basis complex-coefficient models.}
\label{fig:wavepacket-error-en}
\end{figure}

\begin{figure}[!htbp]
\centering
\includegraphics[width=0.88\textwidth]{results/figures/wavepacket_shared_spectral_ltpb/phase_slope_train_midpoint.png}
\caption{The shared network preserves the phase slope at midpoint wavenumbers
that were excluded from training.}
\label{fig:wavepacket-phase-en}
\end{figure}

\begin{figure}[!htbp]
\centering
\includegraphics[width=0.88\textwidth]{results/figures/wavepacket_shared_spectral_ltpb/wavepacket_density_spacetime.png}
\caption{Probability-density evolution and absolute error for the reconstructed
Gaussian wave packet.}
\label{fig:wavepacket-density-en}
\end{figure}

\begin{figure}[!htbp]
\centering
\includegraphics[width=0.88\textwidth]{results/figures/wavepacket_shared_spectral_ltpb/branch_center_width.png}
\caption{The positive- and negative-energy packets propagate in opposite
directions, with center velocities and widths satisfying the preregistered
criteria.}
\label{fig:wavepacket-branch-en}
\end{figure}

The mean group velocities of the positive- and negative-energy packets are
\(+0.3665\) and \(-0.3659\), respectively. The residual root-mean-square at
untrained midpoint wavenumbers is \(0.0017\)--\(0.0033\). The spinor error
increases from \(0.0066\) in the two-wavenumber experiment to \(0.0410\),
indicating that approximation becomes harder as the represented spectral
interval broadens. For every one of the five random seeds, both the spinor and
current errors remain more than one order of magnitude below the best
trainable baseline.

\noindent\textbf{Stage conclusion.}
All preregistered criteria are satisfied. The phase representation has been
extended from discrete energy components to continuous functions
\(S_\pm(k,t)\) over a prescribed interval of the free Dirac spectrum. Momentum
coupling induced by an external potential remains outside the present model.

\section{Synthesis of the current evidence}

\subsection{Established findings}

\begin{enumerate}
  \item \textbf{Scope of a single common phase.}
  It is effective for a pure positive-energy state. For a
  positive/negative-energy superposition, one of the rapid temporal factors
  remains in the complex amplitude and the fixed-budget error increases.
  \item \textbf{Role of eigenspace decomposition.}
  Projection onto the positive- and negative-energy eigenspaces, followed by
  an independent phase for each component, recovers the relative phase and
  current oscillation.
  \item \textbf{Key structural evidence.}
  Parameter-matched ablations associate low error with the joint use of
  eigenspace alignment and phase separation.
  \item \textbf{Validated scope.}
  Tests now cover a single-wavenumber superposition, a two-wavenumber discrete
  spectrum, and a Gaussian packet over a prescribed interval of the free
  spectrum.
\end{enumerate}

\subsection{Accuracy as problem complexity increases}

\begin{table}[htbp]
\centering
\small
\caption{Accuracy across problems of increasing spectral complexity}
\begin{tabularx}{\textwidth}{@{}Y Y r r r@{}}
\toprule
Problem & Method & Spinor error & Current error & Fidelity \\
\midrule
Single-wavenumber positive/negative superposition & LTPB &
0.0025 & 0.0046 & 0.999996 \\
Two wavenumbers, four energy components & Multi-LTPB &
0.0066 & 0.0110 & 0.999965 \\
33-mode Gaussian packet & Shared-Spectral-LTPB &
0.0410 & 0.0481 & 0.998249 \\
\bottomrule
\end{tabularx}
\end{table}

As the number of spectral components increases, the spinor error rises from
\(0.0025\) to \(0.0410\), while fidelity remains above \(0.998\). Phase
separation continues to provide a stable coordinate system for the free
problem, although the remaining accuracy margin narrows with spectral width.
The central next question is whether such coordinates remain stable when
momentum channels become dynamically coupled.

\subsection{Relationship to prior work and prospective advantages}

Kumar et al.\ introduced CPA-PINNs for semiclassical Schr\"odinger-type
equations, combining a WKB-style phase-amplitude decomposition with continuity,
mass, and energy constraints for forward reconstruction and noisy inverse
problems \cite{kumar2026}. In the scalar Schr\"odinger setting, the probability
current explicitly contains the spatial gradient of the common phase. The
Dirac current is the spinor bilinear
\(\psi^\dagger\sigma_x\psi\), from which a global phase cancels. LTPB retains
rapid-phase extraction while assigning phases according to the spectral
projectors of the Dirac Hamiltonian. Its expected advantage therefore concerns
long-time evolution with coherent contributions from multiple energy branches.

Lorin and Yang developed a general PINN formulation for the time-dependent
Dirac equation and studied norm constraints, curved-space settings, and
fermion-doubling issues \cite{lorin2022}. Their direct spinor representation
requires little prior spectral structure. LTPB introduces a stronger physical
coordinate system when the free spectral decomposition is known. Under the
present matched budgets, it reduces the single-mode mixed-state spinor error
from \(0.7991\) for SP-CPA to \(0.0025\), and lowers the 33-mode wave-packet
error by a factor of about 22 relative to the best trainable baseline. These
results support the hypothesis that eigenspace-aligned coordinates improve
optimization conditioning and long-time phase accuracy for spectrally
resolvable, strongly oscillatory Dirac problems.

Fourier pseudospectral and multiscale time-integrator methods provide
high-accuracy forward solvers with rigorous error estimates, including in the
nonrelativistic limit \cite{bao2016}. The current evidence does not establish a
single-query forward-cost advantage for LTPB. Its prospective role lies in
combining governing equations with sparse noisy observations, constructing a
continuously differentiable space-time surrogate, and performing state or
parameter inference in the same computational framework. These uses remain
research hypotheses.

\section{Next stage}

\subsection{Stage VII: weak smooth periodic potential}

The first external-potential test uses
\begin{equation}
V(x)=V_0\cos(qx)I_2.
\end{equation}
In Fourier space, neighboring momentum channels satisfy
\begin{equation}
i\partial_t\widetilde\psi_k=
H(k)\widetilde\psi_k+\frac{V_0}{2}
\left(\widetilde\psi_{k-q}+\widetilde\psi_{k+q}\right).
\end{equation}
This model introduces nontrivial momentum transfer while retaining a smooth
potential and a transparent physical mechanism. The experiment will test
recovery of modal population exchange, current density, and coherence between
energy components. At low-occupancy or newly generated modes, the amplitude
approaches zero and a polar phase can become ill-conditioned. Candidate
extensions include phase regularization near zero amplitude and hybrid
complex/polar coordinates.

\subsection{Subsequent stages}

\begin{itemize}
  \item \textbf{Stage VIII:} add sparse noisy density or current observations
  for state reconstruction and inference of \(m\) and potential parameters.
  \item \textbf{Stage IX:} conduct a unified benchmark against adaptive-sampling
  PINNs, loss-balancing methods, Fourier spectral solvers, and split-operator
  methods.
  \item \textbf{Stage X:} assemble the theory, architecture, ablations,
  applicability, data experiments, complexity, and failure boundaries into a
  submission-ready study.
\end{itemize}

\section{Current scope of the conclusions}

The available experiments support the following statement: for the
one-dimensional free Dirac equation, a prescribed Fourier expansion, and a
fixed training budget, independent learnable phases aligned with the positive-
and negative-energy eigenspaces substantially reduce the spinor, current, and
relative-phase errors of mixed-energy states.

The following questions remain open:
\begin{itemize}
  \item automatic identification of unknown spectra and eigenspaces from data
  and equation residuals;
  \item validity under general external potentials, in higher dimensions, and
  for strongly coupled Dirac equations;
  \item forward accuracy and computational efficiency relative to mature
  spectral and split-operator solvers;
  \item generality, strict novelty, and reproducible boundaries of the method.
\end{itemize}

A submission at the intended journal or conference level requires at least one
external-potential experiment, a sparse noisy-data application, classical
solver references, and a complete cost analysis.

\enlargethispage{3\baselineskip}
\begin{thebibliography}{9}
\raggedright
\footnotesize

\bibitem{kumar2026}
S. Kumar, H. K. Dwivedi, and Rajeev,
\emph{Structure-preserving conservative phase-amplitude PINNs for semiclassical
Schr\"odinger-type equations},
\textit{Journal of Computational and Applied Mathematics}, vol. 485, 117607,
2026.
\href{https://doi.org/10.1016/j.cam.2026.117607}{doi:10.1016/j.cam.2026.117607}.

\bibitem{lorin2022}
E. Lorin and X. Yang,
\emph{Time-dependent Dirac equation with physics-informed neural networks:
Computation and properties},
\textit{Computer Physics Communications}, vol. 280, 108474, 2022.
\href{https://doi.org/10.1016/j.cpc.2022.108474}{doi:10.1016/j.cpc.2022.108474}.

\bibitem{bao2016}
W. Bao, Y. Cai, X. Jia, and Q. Tang,
\emph{A uniformly accurate multiscale time integrator pseudospectral method for
the Dirac equation in the nonrelativistic limit regime},
\textit{SIAM Journal on Numerical Analysis}, vol. 54, no. 3, pp. 1785--1812,
2016.
\href{https://doi.org/10.1137/15M1032375}{doi:10.1137/15M1032375}.

\end{thebibliography}

\section{Project source material}

\begin{itemize}
  \item Research proposal: limitations of CPA-PINN transfer to the Dirac
  equation and a spinor dual-phase formulation.
  \item Direct-transfer report: \path{EXPERIMENT_REPORT.md}
  \item LTPB pilot report: \path{LEARNED_TWO_PHASE_PILOT_REPORT.md}
  \item Mechanism-ablation preregistration:\newline
  \path{LTPB_MECHANISM_ABLATION_PREREGISTRATION.md}
  \item Mechanism-ablation report:
  \path{LTPB_MECHANISM_ABLATION_REPORT.md}
  \item Two-wavenumber preregistration:\newline
  \path{MULTIMODE_LTPB_PREREGISTRATION.md}
  \item Two-wavenumber report:
  \path{MULTIMODE_LTPB_EXPERIMENT_REPORT.md}
  \item Wave-packet preregistration:\newline
  \path{WAVEPACKET_SHARED_SPECTRAL_LTPB_PREREGISTRATION.md}
  \item Wave-packet report:\newline
  \path{WAVEPACKET_SHARED_SPECTRAL_LTPB_REPORT.md}
\end{itemize}

All figures are referenced directly from the existing
\path{results/figures/} directory.

\end{document}
